\documentclass[11pt,reqno]{amsart}
\usepackage[T1]{fontenc}
\usepackage[utf8]{inputenc}
\usepackage{lmodern,amsmath,amssymb,amsthm,mathtools,booktabs,array}
\usepackage[margin=1in]{geometry}
\usepackage{microtype}
\usepackage[hidelinks]{hyperref}
\ifdefined\pdfinfoomitdate\pdfinfoomitdate=1\fi
\ifdefined\pdftrailerid\pdftrailerid{}\fi
\allowdisplaybreaks
\setlength{\emergencystretch}{3em}
\newtheorem{theorem}{Theorem}[section]
\newtheorem{lemma}[theorem]{Lemma}
\newtheorem{proposition}[theorem]{Proposition}
\newtheorem{corollary}[theorem]{Corollary}
\theoremstyle{remark}\newtheorem{remark}[theorem]{Remark}
\newcommand{\R}{\mathbb R}\newcommand{\N}{\mathbb N}
\newcommand{\E}{\mathbb E}\newcommand{\Pp}{\mathbb P}
\newcommand{\Bcal}{\mathfrak B}\newcommand{\KL}{\operatorname{KL}}
\newcommand{\op}{\mathrm{op}}\newcommand{\Jop}{\mathsf J}
\newcommand{\diam}{\operatorname{diam}}
\title[Energy-robust boundary identification]{Boundary Identification from Recurring Fast Probes:\\Energy-Robust Posteriors and Bound Observable Certificates}
\author{Qian Qi}
\date{September 5, 2026}
\begin{document}
\begin{abstract}
We recover logarithmically many labeled physical coefficients of an infinite,
pinned and damped Jacobi lattice from bounded probes without resetting the
state. A recurring sufficiently short duration suffices; the exploration law
may be atomless when nonexploratory short probes have positive probability.
A channel-corrected sampled logarithm has an exact rational finite transform
and a remainder governed by the actual clock norm. Together with a
uniform-in-tail physical inverse and all-sign block energy, it gives an
explicit region $C_\tau r+p_\tau s<1$ for depth $\lfloor r\log n\rfloor$
and error $n^{-s}$. The likelihood uses every completed readout, with bounded
feedback between precommitted blocks and an arbitrary working initial-state
law on a fixed ball. We prove a separate chronological posterior transfer
inequality for predictable readout discrepancies of cumulative squared energy
$D_n$. It retains the region whenever $D_n=o(n^{1-C_\tau r-p_\tau s}/\log n)$,
including non-summable, polynomially decaying discrepancies. We also give a
compatible finite-response estimator, a dimension-uniform finite-section
bound, and an all-policy logarithmic-depth lower bound in readout count.
An executable rational enclosure contract binds the model, clock, complete
observable grid, bands and error budgets before returning a diameter decision.
The elapsed-time order is asserted for experiments with a positive sampling
gap; no unrestricted sampling-rate or sharp joint-frontier claim is made.
\end{abstract}
\maketitle
\tableofcontents

\section{The result and its statistical interfaces}
A short forcing interval and the age at which its response is observed are
different resources. Retained state exposes a short pulse at many later lags.
The target here is a labeled physical prefix, not a realization modulo a
similarity transform. The following theorem gives the experiment precisely.

\begin{theorem}[Recurring fast probes and logarithmic physical depth]
\label{thm:main}
Fix the coefficient box of Section~\ref{sec:model}, a duration cap $T>0$,
noise variance $\sigma^2>0$, probing amplitude $a>0$, baseline bound $U_0$,
and initial-state radius $R_0<\infty$. Choose $0<t_0\le T$ with
$2t_0\Lambda\le1/64$. Let $\chi$ be any probability on $(0,T]$, put
$w=\chi(\{t_0\})$, and choose $\rho\in[0,1]$ such that
\[
 p=1-\rho+\rho w>0.
\]
Each stage consists of a gap of length $t_0$ followed by a probe. Independently
choose its duration from $\chi$ with probability $\rho$, and otherwise use
$t_0$. Add an independent force $aS_i$, $S_i\in\{-1,1\}$ equiprobably,
on every probe. Before a block starts, commit all baseline profiles, bounded
by $U_0$, including their dependence on duration choices, without using any
sign or readout from that block. Feedback from earlier blocks is unrestricted
within this bound. The state is never reset.

For the explicit constants in \eqref{eq:clock-constants}, let $r,s>0$ satisfy
$\gamma=C_\tau r+p_\tau s<1$. There is a block length $m_n=O(\log n)$ such
that, at $J_n=\lfloor r\log n\rfloor$, $\delta_n=n^{-s}$ and
$n'=m_n\lfloor n/m_n\rfloor$, the complete working posterior satisfies
\[
 \Pi_{n'}^{\zeta_{n'}}\{d_{J_n}\ge\delta_n\}\longrightarrow0
 \quad\hbox{in probability}.
\]
The convergence is uniform over admissible truths, bounded external initial
states, product coefficient priors with a fixed positive lower density bound,
working initial-state laws on the radius-$R_0$ ball, and these committed
profiles. Coordinatewise posterior medians recover the prefix to error
$\delta_n$ with probability tending to one. Deterministically,
\[
 t_0n'\le\mathsf T_{n'}\le(T+t_0)n'.
\]
A separate finite-response estimator attains the same sufficient region in
this correctly specified physical experiment. The all-policy bound of
Theorem~\ref{thm:global} matches its fixed-accuracy depth order in readout count.
\end{theorem}

When $\rho<1$, no atom of $\chi$ is required: take $w=0$ and count the
nonexploratory short stages. At $\rho=0$ the duration law is irrelevant.
At $\rho=1$ this construction requires $w>0$. This last requirement is a
condition on the particular information certificate, not an impossibility
statement about other experiments. The chosen fast clock and within-block
commitment are explicit restrictions. All unsuccessful blocks remain in the
likelihood.

\begin{theorem}[Extension to cumulative-energy discrepancies]
\label{thm:main-drift}
In the same experiment, replace the observation equation by
\[
 Y_i=m_i^{\beta_0}+d_i+\xi_i,
\]
where $d_i$ is predictable before the fresh Gaussian noise $\xi_i$ and
$\sum_{i\le n'}d_i^2\le D_{n'}$ pathwise for a deterministic envelope.
The working likelihood is still \eqref{eq:likelihood}; it need not model $d_i$.
The posterior and median conclusions of Theorem~\ref{thm:main} hold if
\begin{equation}
 D_{n'}=o\bigl(n^{1-\gamma}/\log n\bigr).
 \label{eq:main-drift-rate}
\end{equation}
The bound is uniform for fixed uniform little-oh envelopes. The discrepancy
may include the true initial-state response and additional predictable error.
No assertion of an unbiased response estimator under arbitrary sign-correlated
$d_i$ is made.
\end{theorem}

For instance, an additional discrepancy bounded by $C i^{-\alpha}$ with
$0<\alpha\le1$ can fail to be summable. The posterior still attains every
$\gamma<\min(1,2\alpha)$ (Corollary~\ref{cor:polynomial-drift}). Thus the
new robustness statement is not limited to an exponentially transient initial
condition. Its budget is relative to the available information, not to $n$
alone, and it does not permit arbitrary fixed nonvanishing misspecification.

\subsection{Contribution by precise interfaces}
The classical moment and spectral inverse for Jacobi matrices is the
starting point \cite{Teschl,MM2019}, not a new principle claimed here.
The specific quantitative chain is: a channel-corrected pulse transform;
an inverse Lipschitz constant uniform in the infinite tail; a chronological
block-information floor; and a full-posterior transfer with an explicit
squared-energy perturbation budget. The latter transfer is also formulated
as an abstract contrast-family assertion in Theorem~\ref{thm:energy}.
Sections~\ref{sec:transfer} and \ref{sec:comparison} expose what an existing
finite-response estimate can supply, rather than treating latent dimension
or bounded excitation as a general obstruction.

Bayesian robustness to approximated likelihoods is a wider subject: for
example, Seizilles--Siebel \cite{SS2026} study heteroscedastic nonlinear
regression using proxy variances and forward maps. The statement proved here
instead tracks a predictable chronological contrast family, conditional block
information and an arbitrary integrated transient law. We neither import
their theorems nor assert priority for posterior robustness in general.
The explicit positive theorem, rather than an assertion of journal suitability,
is the mathematical contribution submitted for examination.

\section{Model and explicit coefficient inversion}
\label{sec:model}
Fix $0<c_-<c_+$, $0<a_-<a_+$ and $2a_+<b_-<b_+$, all finite, and set
\[
 \Bcal=[c_-,c_+]\times\prod_{k\ge0}[a_-,a_+]\times\prod_{k\ge0}[b_-,b_+].
\]
For $\beta=(c,(a_k),(b_k))$, let $\Jop_\beta$ have diagonal $b_k$ and
negative adjacent entries $-a_k$ on $\ell^2(\N_0)$. Put
\[
 A_\beta=\begin{pmatrix}0&I\\-\Jop_\beta&-cI\end{pmatrix},\quad
 B=\binom0{e_0},\quad\ell(q,v)=q_0,\quad
 g_\beta(t)=\ell e^{tA_\beta}B,\quad h_\beta(t)=\int_0^t g_\beta(u)\,du.
\]
The force $u(t)$ enters through $B$; a state $z$ produces readout
$\ell e^{tA_\beta}z+(g_\beta*u)(t)$. Write
\[
 d_J(\beta,\gamma)=\max\{|c-\widetilde c|,
 |a_k-\widetilde a_k|,|b_k-\widetilde b_k|:0\le k\le J\},
\]
with the evident convention that the damping term is included once. Set
\[
 k_-=b_--2a_+,
 B_0=\max(1,b_++2a_+),\quad \Lambda=1+c_++B_0,\quad
 Q=8\max(2,\Lambda,T,T^{-1},a_-^{-1}).
\]
Then $\|A_\beta\|\le\Lambda$, $Q\ge16$, and
$k_-I\le\Jop_\beta\le B_0I$.

\begin{lemma}[Common stability and compact response geometry]
\label{lem:stability}
There are explicit constants $C_*\ge1$, $\lambda_*>0$ with
$\|e^{tA_\beta}\|\le C_*e^{-\lambda_*t}$ uniformly on the box.
The maps $\beta\mapsto g_\beta$ into $L^1(0,\infty)$ and
$\beta\mapsto h_\beta$ into $C([0,T])$ are continuous and have compact images.
\end{lemma}
\begin{proof}
Let $\epsilon=\min(1,c_-/2,\sqrt{k_-}/2)$ and
\[
 E(q,v)=\tfrac12(\|v\|^2+\langle\Jop q,q\rangle)
       +\epsilon\langle q,v\rangle+\epsilon c\|q\|^2/2.
\]
Set $m_E=\min(1,k_-)/4$ and
$M_E=\max(3/4,B_0/2+\epsilon^2+\epsilon c_+/2)$.
Young's inequality gives $m_E\|z\|^2\le E(z)\le M_E\|z\|^2$ and
$\dot E\le-d_E\|z\|^2$, $d_E=\min(c_-/2,\epsilon k_-)$.
Take $C_*=\sqrt{M_E/m_E}$ and $\lambda_*=d_E/(2M_E)$.
Coordinate convergence implies strong operator convergence, first on finite
support and then by uniform boundedness. Exponential series converge uniformly
on bounded time intervals. Apply $B,\ell$ and then the common integrable
stability tail. Compactness follows from the product compactness of $\Bcal$.
\end{proof}

Here and below vector coefficient norms are one-norms or maximum norms as
specified; matrix infinity norms are induced row-sum norms. Let
$\mu_m=\langle e_0,\Jop^m e_0\rangle$, $v_r=h^{(r+2)}(0)$.
The identity $A^2+cA=-\operatorname{diag}(\Jop,\Jop)$ gives
\begin{equation}
 c=-v_1,\qquad
 \mu_m=(-1)^m\sum_{k=0}^m\binom mk c^{m-k}v_{m+k}.
 \label{eq:moment}
\end{equation}
Indeed expand $(A^2+cA)^m A$ and apply $\ell,B$. For input error
$e=\max_{r\le2M}|v_r-\widetilde v_r|$, $M\ge1$, subtraction yields
\[
 |\mu_m-\widetilde\mu_m|
 \le\{(1+c_+)^m+m2^{m-1}Q^{2m}\}e\le Q^{4m}e\quad(1\le m\le M).
\]
Use $|v_r|\le Q^{r+1}$, $|c-\widetilde c|\le e$, and
$m2^{m-1}\le4^m\le Q^m$; $\mu_0=1$ is fixed.

Let $p_k$ be the monic orthogonal polynomial for the boundary spectral measure,
$\rho_k=\int p_k^2\,d\mu$, and $H_k=(\mu_{r+s})_{r,s=0}^k$.
The Jacobi recurrence gives
\[
 \rho_k=\prod_{r<k}a_r^2,\quad
 p_k(x)=\det(xI-\Jop_{[0,k-1]}),\quad \|p_k\|_1\le(1+B_0)^k\le Q^k.
\]
Stacking padded monic vectors as the rows of a unitriangular matrix gives
\begin{equation}
 H_{k-1}^{-1}=\sum_{r<k}\rho_r^{-1}p_rp_r^T,\qquad
 \|H_{k-1}^{-1}\|_\infty\le\sum_{r<k}Q^{4r}\le Q^{5k}.
 \label{eq:gram}
\end{equation}
The Gram equations for the lower coefficients of $p_k$ have right side
$-(\mu_k,\ldots,\mu_{2k-1})^T$. If the moments have maximum error $e$, set
\[
 S_k=Q^k,\quad A_k=Q^{5k},\quad V_k=kA_k(1+kS_k),\quad l_k=a_-^{2k},
\]
\[
 R_k'=S_k^2+2B_0^{2k}S_kV_k,\qquad T_k'=S_k^2+2B_0^{2k+1}S_kV_k.
\]
Subtracting the equations gives $\|p_k-\widetilde p_k\|_1\le V_ke$.
The errors of $\rho_k=p_k^TH_kp_k$ and
$t_k=\sum p_{k,r}p_{k,s}\mu_{r+s+1}$ are at most $R_k'e,T_k'e$:
one moment change costs $S_k^2e$ and the two polynomial changes give the
other terms. Since $b_k=t_k/\rho_k$ and $a_k^2=\rho_{k+1}/\rho_k$, their
secant constants are
\[
 (T_k'+b_+R_k')/l_k,\qquad
 (R_{k+1}'+a_+^2R_k')/(2a_-l_k).
\]
For $k\ge1$ the successive bounds are
\[
 V_k\le Q^{9k},\quad R_k'\le Q^{13k},\quad T_k'\le Q^{14k},\quad
 \frac{T_k'+b_+R_k'}{l_k}\le Q^{17k},\quad
 \frac{R_{k+1}'+a_+^2R_k'}{2a_-l_k}\le Q^{15(k+1)}.
\]
Before simplification the exponents are at most $8k+1,12k+1,12k+2,16k+1$
and $15k+14$; retain the factor two in the last ratio. At $k=0$, $b_0=\mu_1$
has constant one and $a_0$ has constant at most $Q^{14}/2$ because $\rho_0=1$.
Thus moments through $2J+2$ reconstruct the coefficients through $J$ with
constant $Q^{17(J+1)}$.

\begin{theorem}[All-depth jet inverse]
\label{thm:jet}
Put $j=J+1$, $R=4J+8$ and $L=Q^{25j}$. For all admissible pairs,
\begin{equation}
 d_J(\beta,\gamma)\le L\max_{0\le r\le R}
 |h_\beta^{(r)}(0)-h_\gamma^{(r)}(0)|.
 \label{eq:jet}
\end{equation}
\end{theorem}
\begin{proof}
Take $M=2J+2$ in \eqref{eq:moment}. The two constants have exponents
$8j$ and $17j$. Derivatives through $4J+6$ suffice; the two spare orders
preserve the earlier diagnostic law. The calculation includes $J=0$ and uses
only finite moment identities, uniformly in the unspecified tail.
\end{proof}
Bounded-generator analyticity and \eqref{eq:jet} prove injectivity of the response.
The corresponding classical Weyl relation is
$\widehat h(z)=z^{-1}m_0(z^2+cz)$ and
$m_k(w)^{-1}=w+b_k-a_k^2m_{k+1}(w)$ \cite{Teschl}.


\section{Clock-sensitive sampled inversion}
\label{sec:sampled}
Put $\Delta=2t_0$, $U_\beta=e^{\Delta A_\beta}$,
$X_\beta=U_\beta-I$ and $H_\beta=\int_0^{t_0}e^{uA_\beta}\,du$. Write
\[
 b_{k,\beta}=\ell U_\beta^kH_\beta B
   =h_\beta(k\Delta+t_0)-h_\beta(k\Delta).
\]
The extra subscript separates this pulse response from a physical diagonal
coefficient. Set $x=\Delta\Lambda$ and
\begin{equation}
 \tau=\frac{2x}{1-x}\le\frac2{63}<\frac1{16}.
 \label{eq:tau}
\end{equation}

\begin{lemma}[Pulse transform and actual-clock remainder]
\label{lem:sampled}
Let $N\ge R=4J+8$. There are parameter-independent rational coefficients
$T_{rl}$ at rational clock inputs such that, with
\begin{equation}
 \mathcal A=8\Delta^{-R}4^N,\qquad
 \mathcal E=\Delta^{-R}\tau^N,\qquad
 \mathcal W=\max_{1\le r\le R}\sum_{l=0}^N\frac{T_{rl}^2}{N+1-l},
 \label{eq:transform-budgets}
\end{equation}
we have $0<\mathcal W\le\mathcal A^2$. For every admissible pair, put
$\mathcal D_R=\max_{0\le r\le R}|h_\beta^{(r)}(0)-h_\eta^{(r)}(0)|$.
Then
\begin{align}
 \mathcal D_R&\le\mathcal A\max_{l\le N}|b_{l,\beta}-b_{l,\eta}|+\mathcal E,
 \label{eq:sampled-infty}\\
 \mathcal D_R&\le\sqrt{\mathcal W}
 \left(\sum_{l=0}^N(N+1-l)|b_{l,\beta}-b_{l,\eta}|^2\right)^{1/2}
 +\mathcal E.\label{eq:sampled-weighted}
\end{align}
\end{lemma}
\begin{proof}
The elementary series bound $e^x-1\le x/(1-x)$ gives
$\|X_\beta\|\le x/(1-x)$. In this common norm ball the logarithm and square
root converge absolutely, so $\log(I+X_\beta)=\Delta A_\beta$ and
$(I+X_\beta)^{1/2}=e^{\Delta A_\beta/2}$. Define
\[
 F_r(z)=\Delta^{-r}[\log(1+z)]^{r-1}
 \frac{\log(1+z)}{z}(\sqrt{1+z}+1)=\sum_{k\ge0}c_{rk}z^k.
\]
Analytic continuation defines the removable quotient at zero. Multiplication
by the integrated channel gives $F_r(X_\beta)H_\beta=A_\beta^{r-1}$.
On $|z|=1/2$, the logarithm has modulus at most $\log2<1$, its quotient
has modulus at most $2$, and $|\sqrt{1+z}+1|<3$. Since $\Delta<1$,
$|c_{rk}|\le6\Delta^{-R}2^k$ and $\|H_\beta\|\le\Delta$.
The difference of two degree-$N$ tails is at most
\[
 12\Delta^{1-R}\frac{\tau^{N+1}}{1-\tau}
 =\Delta^{-R}\tau^N\frac{12\Delta\tau}{1-\tau}
 \le\mathcal E,
\]
since the final factor is at most $24\Delta/61\le1$.
Expanding $(U-I)^k$ gives
$T_{rl}=\sum_{k=l}^N c_{rk}(-1)^{k-l}\binom{k}{l}$ and row one-norm at
most $6\Delta^{-R}\sum_{k=0}^N4^k\le\mathcal A$.
Weighted Cauchy--Schwarz proves the second assertion and
$\mathcal W\le\mathcal A^2$. The first row is nonzero: its polynomial at
$U=1$ is $c_{10}=2/\Delta$. Rational series operations prove rationality.
The zeroth derivative difference is zero.
\end{proof}

Let $j=J+1$, $d_0=\log\Delta^{-1}$, $v=\log(1/\delta)$ and
$\ell_\tau=\log(1/\tau)$. For $0<\delta\le1$, choose
\begin{equation}
 N=\max\left\{R,\left\lceil
 \frac{\log2+25j\log Q+Rd_0+v}{\ell_\tau}\right\rceil\right\},\qquad m=N+1.
 \label{eq:choose-N}
\end{equation}
This ensures $L\mathcal E\le\delta/2$. Define
\begin{equation}
 \kappa=p^m\frac{\delta^2}{4L^2\mathcal A^2},\qquad
 \kappa_W=p^m\frac{\delta^2}{4L^2\mathcal W}\ge\kappa.
 \label{eq:kappa}
\end{equation}
The following closed-form constants include the entire block probability:
\begin{equation}
 \begin{split}
 K_\tau&=12+\frac{\log2+25\log Q+8d_0}{\ell_\tau},\qquad
 H_p=\log(1/p)+\log16,\\
 C_\tau&=\log256+50\log Q+16d_0+K_\tau H_p,\qquad
 p_\tau=2+H_p/\ell_\tau.
 \end{split}\label{eq:clock-constants}
\end{equation}

\begin{proposition}[Resource envelope]
\label{prop:envelope}
For all $J\ge0$ and $0<\delta\le1$,
$m\le K_\tau j+v/\ell_\tau$ and
$\kappa_W\ge\kappa\ge e^{-C_\tau j}\delta^{p_\tau}$.
\end{proposition}
\begin{proof}
Use $R\le8j$, $j\ge1$ and a ceiling error at most one in
\eqref{eq:choose-N}. Directly,
\[
 -\log\kappa=m\log(1/p)+2v+\log256+50j\log Q+2Rd_0+N\log16
 \le C_\tau j+p_\tau v.
\]
\end{proof}
For fixed $p>0$, smaller chosen clocks can make $p_\tau$ approach $2$.
The depth constant changes at the same time and can worsen. This is an
accuracy-exponent refinement, not a sharp fixed-depth tradeoff. The exact
weighted certificate may further improve a finite instance without changing
the conservative closed-form region.

\section{The complete working likelihood and variance-sensitive transfer}
\label{sec:likelihood}
Let $U_*=U_0+a$. Each readout noise $\xi_i$ is a fresh $N(0,\sigma^2)$
variable, observed after the current actions. All policy kernels are
parameter-independent functions of the recorded past. Write $m_i^\beta$ for
the zero-state response to the \emph{entire} realized input,
$b_i^\beta(z)=\ell e^{\mathsf T_iA_\beta}z$, and
\begin{equation}
 L_n^{\zeta_n}(\beta)=\int\prod_{i\le n}
 \varphi_\sigma(Y_i-m_i^\beta-b_i^\beta(z))\,\zeta_n(dz),\qquad
 \Pi_n^{\zeta_n}(d\beta)\propto L_n^{\zeta_n}(\beta)\Pi(d\beta).
 \label{eq:likelihood}
\end{equation}
Here $\zeta_n$ is any probability on $\|z\|\le R_0$; its support need not
contain the true external initial state. This is the exact \emph{working}
likelihood, not a correct-specification assertion. Action kernels cancel at
a common history in a chronological likelihood ratio; one never conditions
on the terminal, response-dependent design as though it were fixed.

Let $f_i(\beta)=m_i^\beta-m_i^{\beta_0}$ and put
\[
 d(\beta,\gamma)=(U_*+a)\|g_\beta-g_\gamma\|_1,\qquad
 M=\max(1,2U_*C_*/\lambda_*).
\]
The compact metric space follows from Lemma~\ref{lem:stability} and
Theorem~\ref{thm:jet}. Uniformly over all histories,
$|f_i|\le M$ and $|f_i(\beta)-f_i(\gamma)|\le d(\beta,\gamma)$.
Since $\mathsf T_i\ge t_0i$, the initial-state envelope is
$r_i=C_*R_0e^{-\lambda_*t_0i}$. Define
\[
 R_1=\frac{C_*R_0}{e^{\lambda_*t_0}-1},\quad
 R_2=\frac{C_*^2R_0^2}{e^{2\lambda_*t_0}-1},\quad
 b_0=3(MR_1+R_2)/\sigma^2,\quad b_1=2R_1/\sigma.
\]

\subsection{Deterministic covers and prior balls}
Suppose every coordinate density of the product prior is at least $p_->0$.
Let $w_->0$ be the smallest coordinate interval width and let $W\ge1$ bound
the largest width. Put $C_d=\max(1,4(U_*+a)C_*^2/\lambda_*^2)$.
For $0<\epsilon<1$ set
\begin{align}
 S(\epsilon)&=\max\{1,\lambda_*^{-1}
 \log(16(U_*+a)C_*/(\lambda_*\epsilon))\},\notag\\
 K(\epsilon)&=\min\{K\ge0:2(U_*+a)S e^{\Lambda S}
 (\Lambda S)^{K+1}/(K+1)!\le\epsilon/8\},\notag\\
 r(\epsilon)&=\min(w_-,1,\epsilon/(4C_d)),\notag\\
 \overline{\mathcal H}(\epsilon)&=(2K+3)\log(2+W/r(\epsilon)),\qquad
 \mathcal P(\epsilon)=(2K+3)\log(1/(p_-r(\epsilon))).
 \label{eq:entropy}
\end{align}
For $\epsilon\ge1$ use the bounds at $1/2$.

\begin{lemma}[Entropy and thickness]
\label{lem:entropy}
The log covering number in $d$ is at most $\overline{\mathcal H}(\epsilon)$;
every open $d$-ball of radius $\epsilon$ has prior mass at least
$e^{-\mathcal P(\epsilon)}$. Both bounds are $O((1+|\log\epsilon|)^2)$.
\end{lemma}
\begin{proof}
Agreement through $c,a_0,b_0,\ldots,a_K,b_K$ implies agreement of boundary
powers through order $K$. Local exponential series and the common stable tail give
\[
 \|g_\beta-g_\gamma\|_1\le
 2S e^{\Lambda S}(\Lambda S)^{K+1}/(K+1)!+
 2C_*e^{-\lambda_*S}/\lambda_*.
\]
This costs at most $\epsilon/4$ in $d$. A prefix coordinate change of maximum
size $r$ with the same tail changes the generator by at most $4r$; stable
Duhamel factors bound its $L^1$ response change by $4rC_*^2/\lambda_*^2$.
Insert an admissible hybrid to combine the estimates. A finite prefix mesh
of radius $r(\epsilon)$ with a fixed admissible tail is an $\epsilon/2$ cover
with at most $(2+W/r)^{2K+3}$ points. Each coordinate interval about a truth,
including a boundary truth, has length at least $r$; the unrestricted tail has
prior mass one. This proves the mass bound. Finally $S=O(1+|\log\epsilon|)$,
and the factorial condition follows by taking $K+1$ at least
\[
 \max\{2e\Lambda S,
 [\log(16(U_*+a)S/\epsilon)+\Lambda S]/\log2,1\}.
\]
Thus $K=O(1+|\log\epsilon|)$.
\end{proof}

Partition the $n$ readouts into consecutive deterministic groups $G_b$, each
of size at most $m$, and let $\mathcal F_{b,-}$ be its starting history. Set
\[
 A_n=\sum_{i\le n}f_i^2,\qquad S_n=\sum_{i\le n}\xi_i f_i,\qquad
 V_n=\sum_b\E[\sum_{i\in G_b}f_i^2\mid\mathcal F_{b,-}].
\]
These are functions of $\beta$; the groups need not be independent. Define
\begin{equation}
 \begin{split}
 z&=3\sigma+1,\quad e_q=\min\{1,q/[256(M+z)]\},\quad
 C_m=1024\max\{1,\sigma^2,mM^2\},\\
 \Phi_{n,m}(q)&=3\exp\{\overline{\mathcal H}(e_q)-nq/C_m\}+e^{-n}.
 \end{split}\label{eq:variance-tail}
\end{equation}

\begin{theorem}[Score and grouped-square control]
\label{thm:variance}
Except with probability at most $\Phi_{n,m}(q)$, simultaneously for all $\beta$,
\begin{equation}
 |S_n|\le A_n/16+nq/64,\qquad A_n\ge V_n/2-nq/32.
 \label{eq:variance-event}
\end{equation}
\end{theorem}
\begin{proof}
At a deterministic net point, Gaussian chronology makes
$\exp\{\lambda S_n-\lambda^2\sigma^2A_n/2\}$ a nonnegative martingale.
Using $\lambda=1/(8\sigma^2)$ with either sign yields
$\Pp\{|S_n|>A_n/16+8\sigma^2x\}\le2e^{-x}$.
For $X_b=\sum_{i\in G_b}f_i^2\in[0,mM^2]$ and $\alpha_0=1-e^{-1}$,
the chord bound $e^{-u}\le1-\alpha_0u$ on $[0,1]$ gives
\[
 \E\left[\exp\left\{-\frac{X_b}{mM^2}
 +\frac{\alpha_0\E[X_b\mid\mathcal F_{b,-}]}{mM^2}\right\}
 \middle|\mathcal F_{b,-}\right]\le1.
\]
Hence $\Pp\{A_n<\alpha_0 V_n-mM^2x\}\le e^{-x}$.
Choose $x_s=nq/(1024\sigma^2)$ and $x_v=nq/(64mM^2)$.
The event $n^{-1}\sum|\xi_i|\le z$ fails with probability at most $e^{-n}$,
using $\E e^{|\xi_i|/\sigma}\le2e^{1/2}$.
On an $e_q$ net, score interpolation costs at most
$ne_qz+ne_qM/8\le nq/256$; adding $8\sigma^2x_s=nq/128$ is less than
$nq/64$. Interpolating $A_n-\alpha_0V_n$, including its conditional
expectations, costs at most $4Mne_q\le nq/64$. Adding $mM^2x_v=nq/64$
and using $\alpha_0>1/2$ proves the second bound. Union-bound the deterministic
net. Compact continuous fields have measurable suprema on a countable dense
set, so no terminal-design conditioning or independence of groups is used.
\end{proof}

\begin{lemma}[Uniform initial-state integral]
\label{lem:nuisance}
The exact working log ratio is
\[
 \log\frac{L_n^{\zeta_n}(\beta)}{L_n^{\zeta_n}(\beta_0)}
 =\sigma^{-2}(S_n-A_n/2)+\mathcal R_n(\beta),\qquad
 \sup_\beta|\mathcal R_n(\beta)|\le B_n,
\]
where $B_n=\sigma^{-2}(2\sum r_i|\xi_i|+3MR_1+3R_2)$ and
$\Pp\{B_n>b_0+b_1\sqrt{2x}\}\le2e^{-x}$ for $x>0$.
\end{lemma}
\begin{proof}
Completing Gaussian squares relative to the zero-state likelihood gives an
integral of
$\exp\{\sigma^{-2}\sum[b_i^\beta(z)(Y_i-m_i^\beta)-b_i^\beta(z)^2/2]\}$.
Its exponent has absolute value at most
$\sigma^{-2}(\sum r_i|\xi_i|+MR_1+3R_2/2)$ uniformly in $z$ and $\beta$.
The log of an integral against a probability lies between the exponent's
infimum and supremum. Apply this at candidate and truth, and bound the true
transient cross term by $MR_1/\sigma^2$.
Jensen with weights $r_i/R_1$ gives
$\E e^{\lambda\sum r_i|\xi_i|}\le2e^{\lambda^2\sigma^2R_1^2/2}$.
Optimize exponential Markov; for $R_1=0$ the remainder vanishes.
\end{proof}

\begin{theorem}[Complete-posterior information transfer]
\label{thm:posterior}
For a measurable set $F$ with $\inf_F V_n\ge nq>0$ pathwise,
\begin{equation}
 \Pi_n^{\zeta_n}(F)\le
 \exp\{\mathcal P(\sqrt q/4)-nq/(16\sigma^2)\}
 \label{eq:posterior-bound}
\end{equation}
except with probability at most
$\Phi_{n,m}(q)+\Pp\{B_n>nq/(64\sigma^2)\}$.
\end{theorem}
\begin{proof}
On \eqref{eq:variance-event}, $A_n\ge15nq/32$ on $F$, so its numerator
log ratio is at most $-97nq/(512\sigma^2)+B_n$.
On the $d$-ball of radius $\sqrt q/4$, $A_n\le nq/16$ deterministically,
and the log ratio is at least $-13nq/(256\sigma^2)-B_n$.
Integrating that ball gives the denominator lower bound from
Lemma~\ref{lem:entropy}. The remainder condition leaves coefficient
$-97/512+13/256+1/32=-55/512\le-1/16$, proving the claim.
\end{proof}

For posterior error $\epsilon_p\in(0,1)$ and failure $\alpha\in(0,1)$,
a sufficient readout budget is the next multiple of $m$ above
\begin{equation}
 \left\lceil\max\left\{\begin{aligned}
 &1,\quad\frac{C_m}{q}
 [\overline{\mathcal H}(e_q)+\log(9/\alpha)],\quad\log(3/\alpha),\\
 &\frac{64\sigma^2}{q}[b_0+b_1\sqrt{2\log(6/\alpha)}],\\
 &\frac{16\sigma^2}{q}[\mathcal P(\sqrt q/4)+\log(1/\epsilon_p)]
 \end{aligned}\right\}\right\rceil.
 \label{eq:posterior-budget}
\end{equation}
Every probability estimate is uniform in the declared model and policy class.


\section{A chronological posterior principle with an energy budget}
\label{sec:energy}
This section separates the nuisance integrated by the working likelihood from
the discrepancy in the true conditional mean. It also records exactly which
part of the argument is independent of the Jacobi model.

\begin{theorem}[Predictable energy-robust posterior transfer]
\label{thm:energy}
Consider a compact metric parameter space and a predictable mean family
$m_i^\beta$, continuous in the parameter at each fixed history, with
contrast $f_i=m_i^\beta-m_i^{\beta_0}$, $|f_i|\le M$, and a deterministic
metric dominating contrast differences uniformly over histories. Suppose
its covering and prior-ball costs are bounded by
$\overline{\mathcal H}$ and $\mathcal P$ as in Lemma~\ref{lem:entropy}.
Use consecutive deterministic groups of size at most $m$, and define
$A_n,S_n,V_n$ as in Section~\ref{sec:likelihood}. Observe
$Y_i=m_i^{\beta_0}+d_i+\xi_i$, with $d_i$ predictable and fresh
$\xi_i\sim N(0,\sigma^2)$. Assume $\sum_{i\le n}d_i^2\le D$ pathwise.

Use a working likelihood integrating shifts $b_i^\beta(z)$ against any
probability law, with the uniform envelope $|b_i^\beta(z)|\le r_i$ and
$R_1=\sum_i r_i<\infty$, $R_2=\sum_i r_i^2<\infty$. The same law is used
at candidate and truth; no support-at-the-truth condition is required.
For a measurable $F$ with $\inf_F V_n\ge nq>0$ pathwise, put
\begin{equation}
 B_{n,D}=\sigma^{-2}\left(2\sum_{i\le n}r_i|\xi_i|
       +2\sqrt{R_2D}+2MR_1+R_2\right).
 \label{eq:energy-B}
\end{equation}
Except on an event of probability at most $\Phi_{n,m}(q)$ from
\eqref{eq:variance-tail},
\begin{equation}
 \Pi_n(F)\le\exp\left\{\mathcal P(\sqrt q/4)
 -\frac{27nq}{256\sigma^2}+\frac{8D}{\sigma^2}+2B_{n,D}\right\}.
 \label{eq:energy-posterior}
\end{equation}
In particular, if $D\le nq/1024$, then
$\Pi_n(F)\le\exp\{\mathcal P(\sqrt q/4)-nq/(16\sigma^2)\}$
except with probability at most
$\Phi_{n,m}(q)+\Pp\{B_{n,D}>nq/(64\sigma^2)\}$.
For $u>0$, with probability at least $1-2e^{-u}$,
\begin{equation}
 B_{n,D}\le\sigma^{-2}
 \left(2\sigma R_1\sqrt{2u}+2\sqrt{R_2D}+2MR_1+R_2\right).
 \label{eq:energy-tail}
\end{equation}
\end{theorem}
\begin{proof}
The proof of Theorem~\ref{thm:variance} uses only predictability, Gaussian
fresh noise, the bounded group increments, and the deterministic cover.
Consequently its event remains valid under the present observation law;
no correct-specification premise enters that martingale calculation.

Let $L_n^0$ be the zero-shift Gaussian likelihood. Completing squares gives
\[
 \log\frac{L_n^0(\beta)}{L_n^0(\beta_0)}
 =\sigma^{-2}\left(S_n-A_n/2+\sum_i d_i f_i\right).
\]
For a candidate, the exponent in $L_n(\beta)/L_n^0(\beta)$ is
$\sigma^{-2}\sum_i[b_i^\beta(z)(\xi_i+d_i-f_i)
 -(b_i^\beta(z))^2/2]$. Its absolute value is at most
\[
 \sigma^{-2}\left(\sum_i r_i|\xi_i|+\sqrt{R_2D}+MR_1+R_2/2\right)
\]
uniformly in $\beta,z$, by Cauchy--Schwarz. The log of its integral lies
between its infimum and supremum. Subtracting the two log integrals thus
has absolute value at most $B_{n,D}$, before any choice of working law.

Young's inequality gives
$|\sum_i d_i f_i|\le A_n/16+4D$.
On the simultaneous score and square event,
\begin{align*}
 \log\frac{L_n(\beta)}{L_n(\beta_0)}
 &\le\sigma^{-2}(-3A_n/8+nq/64+4D)+B_{n,D},\\
 \log\frac{L_n(\beta)}{L_n(\beta_0)}
 &\ge\sigma^{-2}(-5A_n/8-nq/64-4D)-B_{n,D}.
\end{align*}
On $F$, $A_n\ge15nq/32$, so the numerator coefficient is $-41/256$.
On the prior ball of radius $\sqrt q/4$, $A_n\le nq/16$ pathwise, so the
denominator coefficient is $-7/128$. Its prior mass is at least
$e^{-\mathcal P(\sqrt q/4)}$. Taking the ratio proves
\eqref{eq:energy-posterior}, since $-41/256+7/128=-27/256$.
The particular budgets yield
\[
 -27/256+8/1024+2/64=-17/256\le-1/16.
\]
Finally, for $R_1>0$, Jensen with weights $r_i/R_1$ and
$\E e^{\lambda|\xi_i|}\le2e^{\lambda^2\sigma^2/2}$ gives
$\Pp\{\sum r_i|\xi_i|>\sigma R_1\sqrt{2u}\}\le2e^{-u}$ by exponential
Markov. The $R_1=0$ case is immediate. This proves \eqref{eq:energy-tail}.
\end{proof}

The finite guarantee can be used without an asymptotic assertion: verify
$D\le nq/1024$, the deterministic right side of \eqref{eq:energy-tail}
is at most $nq/(64\sigma^2)$ for $u=\log(4/\alpha)$,
$\Phi_{n,m}(q)\le\alpha/2$, and
$\mathcal P(\sqrt q/4)+\log(1/\epsilon_p)\le nq/(16\sigma^2)$.
These imply posterior error at most $\epsilon_p$ with failure at most $\alpha$.
The square-energy loss $8D/\sigma^2$ is explicit even outside that convenient
sufficient budget; no unmodeled persistent discrepancy is silently removed.

\begin{corollary}[Non-summable polynomial discrepancies]
\label{cor:polynomial-drift}
In Theorem~\ref{thm:main-drift}, suppose the discrepancy is a bounded true
initial-state response plus $e_i$, where $|e_i|\le C i^{-\alpha}$ pathwise
and $\alpha>0$. Every $\gamma<\min(1,2\alpha)$ is attainable for the
complete posterior and its prefix medians. This includes $0<\alpha\le1/2$,
for which the envelope need not have finite total squared energy.
\end{corollary}
\begin{proof}
The initial response has bounded summed squares. The inequality
$(u+v)^2\le2u^2+2v^2$ reduces the envelope to a constant plus
$2C^2\sum_{i\le n}i^{-2\alpha}$. This is $O(n^{1-2\alpha})$ when
$\alpha<1/2$, $O(\log n)$ when $\alpha=1/2$, and $O(1)$ otherwise.
Each is $o(n^{1-\gamma}/\log n)$ under the strict displayed condition.
\end{proof}

\section{Retained-state block information and contraction}
\label{sec:blocks}
A successful block has all $m$ durations equal to $t_0$. Its probability is
$p^m$, generated by exogenous duration choices independently of signs and
readouts. If only a designated atom or nonexploration label is counted,
coincident durations can only enlarge the usable event. Duration randomizers
can be exposed for the energy calculation without conditioning the actual
likelihood on a terminal, response-dependent design.

\begin{lemma}[All-sign information]
\label{lem:block-floor}
Conditional on a successful block's starting history and committed profiles,
put $d_k^{\beta}=b_{k,\beta}-b_{k,\beta_0}$. Then
\begin{equation}
 \E_S\sum_{k=0}^{m-1}f_{b,k}(\beta)^2
 =\sum_{k=0}^{m-1}P_{b,k}(\beta)^2
   +a^2\sum_{k=0}^{m-1}(m-k)(d_k^{\beta})^2.
 \label{eq:all-signs}
\end{equation}
Consequently the conditional group energy is at least $a^2\kappa_W$ on
$\{d_J(\beta,\beta_0)\ge\delta\}$, also under the discrepancies of
Theorem~\ref{thm:main-drift}.
\end{lemma}
\begin{proof}
Along the common realized input,
$f_{b,k}=P_{b,k}+a\sum_{j=0}^k S_{b,j}d_{k-j}^{\beta}$.
Commitment makes the intercept independent of every current block sign.
Sign orthogonality removes cross terms, and each squared response appears
$m-k$ times. Multiply its nonnegative contribution by $p^m$; other duration
patterns contribute nonnegative energy. The jet and weighted sampled inverses,
with $L\mathcal E\le\delta/2$, imply a weighted squared-response energy at
least $\delta^2/(4L^2\mathcal W)$ for separated parameters.
An additive readout discrepancy changes histories and later committed
profiles, but not this conditional algebra along a fixed history. Fresh
signs remain independent of that history. Thus the same floor holds.
\end{proof}

\begin{theorem}[Finite information budget and proofs of the main results]
\label{thm:block-posterior}
At $n'=Bm$, the complete-posterior theorem applies with
$q=a^2\kappa_W/m$ or its simpler value $q=a^2\kappa/m$.
Theorem~\ref{thm:posterior} and budget \eqref{eq:posterior-budget} handle the
bounded true initial state. Theorem~\ref{thm:energy} handles additional
predictable discrepancies without altering the working likelihood.
\end{theorem}
\begin{proof}
Summing Lemma~\ref{lem:block-floor} gives $V_{n'}\ge n'q$ on the entire
separated set. For $J_n=\lfloor r\log n\rfloor$, $\delta_n=n^{-s}$,
Proposition~\ref{prop:envelope} gives $m_n=O(\log n)$,
$\kappa\ge e^{-C_\tau}n^{-\gamma}$ and $n'/n\to1$. Thus
\[
 n'q\ge c\,n^{1-\gamma}/\log n,\qquad
 n'q/C_{m_n}\ge c\,n^{1-\gamma}/(\log n)^2.
\]
Also $\log(1/q)=O(\log n)$ directly from its formula. These quantities
dominate the $O((\log n)^2)$ covering and prior costs. The bounded
initial-state nuisance threshold diverges, proving Theorem~\ref{thm:main}.

For Theorem~\ref{thm:main-drift}, condition \eqref{eq:main-drift-rate} gives
$D_{n'}=o(n'q)$. In \eqref{eq:energy-B} the random weighted-noise term is
uniformly $O_{\Pp}(1)$, while $\sqrt{R_2D_{n'}}=o(n'q)$.
Therefore $B_{n',D_{n'}}=o_{\Pp}(n'q)$, and the same exponent dominates all
costs. The finite-energy theorem proves contraction. Once the posterior has
more than half its mass in $d_J<\delta$, every coordinate median lies within
$\delta$ of its true coordinate; an admissible completion exists in the
product box. Finally every stage has length $t_0+t_i$ between $t_0$ and
$t_0+T$. This proves the elapsed-time bound for the attaining experiment.
\end{proof}

\section{Finite-response estimation and comparison}
\label{sec:transfer}
Let $\mathbf b_N(\beta)=(b_{0,\beta},\ldots,b_{N,\beta})$ with maximum
norm. Its image is compact by Lemma~\ref{lem:stability}.

\begin{theorem}[Response-to-physical transfer]
\label{thm:transfer}
Suppose $\|\widehat{\mathbf b}-\mathbf b_N(\beta_0)\|_\infty\le e$ on
an event of probability at least $1-\alpha$. Take a finite deterministic
$\eta$-net of the response image, assigning an admissible preimage and a
fixed ordering to every point. Minimize the residual over those preimages,
breaking ties by that ordering. Then, on the same event,
\[
 d_J(\widehat\beta,\beta_0)\le L\{\mathcal A(2e+\eta)+\mathcal E\}.
\]
The response confidence set of radius $e$ is nonempty and has $d_J$ diameter
at most $L(2\mathcal A e+\mathcal E)$. Thus $e=\eta=\delta/(6L\mathcal A)$
suffices for estimation, and $e\le\delta/(4L\mathcal A)$ for diameter.
\end{theorem}
\begin{proof}
A net point within $\eta$ of the true response has residual at most $e+\eta$.
The minimizing residual is no larger. The triangle inequality gives response
distance at most $2e+\eta$ from truth. Apply the sampled and jet inverses.
For two members of the confidence set, their response distance is at most
$2e$. The remainder choice completes the claims; a finite ordered minimization
is measurable.
\end{proof}

\begin{theorem}[A compatible nonreset response estimator]
\label{thm:response-estimator}
Under Theorem~\ref{thm:main}'s observation model without an additional
discrepancy, let $M_B$ count successes among $B$ blocks. For $M_B>0$, put
\[
 \widehat b_k=\frac1{aM_B}\sum_{b\le B:\,E_b}S_{b,0}Y_{b,k},\quad
 B_*=C_*(R_0+U_*/\lambda_*),\quad v_*=\sigma^2+B_*^2.
\]
Assign any output at $M_B=0$. For $e>0$,
\begin{equation}
 \Pp\{M_B<Bp^m/2\ \hbox{or}\
 \|\widehat{\mathbf b}-\mathbf b_N(\beta_0)\|_\infty>e\}
 \le e^{-Bp^m/8}+2m e^{-a^2e^2Bp^m/(4v_*)}.
 \label{eq:response-tail}
\end{equation}
This estimator attains the region of Theorem~\ref{thm:main}.
\end{theorem}
\begin{proof}
On success, $Y_{b,k}=aS_{b,0}b_{k,\beta_0}+P^0_{b,k}+\xi_{b,k}$.
Conditional on starting history, success and the other signs, the intercept
is independent of $S_{b,0}$ and bounded by $B_*$. Removing the first signed
pulse leaves an input still bounded by $U_*$. Hence the centered product
$S_{b,0}Y_{b,k}-ab_{k,\beta_0}$ is sub-Gaussian with proxy $v_*$.
For each lag,
\[
 \exp\left\{\lambda\sum_{b\le B}\mathbf1_{E_b}
 (S_{b,0}Y_{b,k}-ab_{k,\beta_0})-\lambda^2v_*M_B/2\right\}
\]
is a nonnegative block supermartingale. Use both signs of $\lambda=ae/v_*$
on $M_B\ge Bp^m/2$, and use the Chernoff count bound $e^{-Bp^m/8}$.
Union over lags proves \eqref{eq:response-tail}; the lag estimates need not
be independent. No conditional Gaussianity at a terminal success count is
asserted. At $e=\eta=\delta/(6L\mathcal A)$, $p^me^2=\kappa/9$.
Since $n\kappa/m$ dominates $\log n$ in the stated region, both failure
terms vanish. The posterior, unlike this auxiliary estimator, uses every block.
\end{proof}

\begin{proposition}[Dimension-uniform approximation]
\label{prop:finite-section}
Cut the chain at site $K\ge1$, retaining the same interior coefficients.
Let $S=N\Delta+t_0$. The cut pulse responses satisfy
\begin{align}
 \max_{k\le N}|b_{k,\beta}-b^{[K]}_{k,\beta}|
 &\le 2t_0e^{\Lambda S}(\Lambda S)^{K+1}/(K+1)!=:\epsilon_K,
 \label{eq:finite-bias}\\
 \sum_{k\ge0}|b^{[K]}_{k,\beta}|
 &\le\frac{C_*t_0e^{\Lambda t_0}}{1-e^{-\lambda_*\Delta}}.
 \label{eq:finite-uniform}
\end{align}
Both bounds are uniform in the unspecified tail and cut size. A sufficient
condition for $\epsilon_K\le\epsilon$ is
\[
 K+1\ge\max\left\{2,2e\Lambda S,
 \frac{\log(2t_0/\epsilon)+\Lambda S}{\log2}\right\}.
\]
In particular $K=O(\log n)$ suffices at logarithmic lag horizon and polynomial
response tolerance.
\end{proposition}
\begin{proof}
Boundary generator words agree through order $K$. Their exponential-series
difference is bounded by $2e^{\Lambda S}(\Lambda S)^{K+1}/(K+1)!$.
Integrate over the short pulse. Principal cuts have the same positive form
bounds and hence the same energy-based stability constants. Apply stability
to $U_K^kH_KB$ and sum the geometric series to obtain
\eqref{eq:finite-uniform}. The last assertion follows from
$q!\ge(q/e)^q$ and $(\Lambda S)^q/q!\le2^{-q}$ for $q\ge2e\Lambda S$.
\end{proof}

\subsection{A theorem-level comparison with finite-time LTI identification}
\label{sec:comparison}
Sarkar--Rakhlin--Dahleh \cite{SRD2020} genuinely allow isotropic sub-Gaussian
inputs in Assumption 1 and in the Section 11 error proof. Their Algorithm 1
uses Gaussian inputs, but that pseudocode is not the input scope of the
analysis. A scalar Rademacher sign has mean zero, variance one and
$\E e^{\lambda S}=\cosh\lambda\le e^{\lambda^2/2}$. It is therefore a
bounded input satisfying the relevant sub-Gaussian requirement.

To compare the experiments, specialize to regular durations and zero baseline.
Let $z_i$ be the state at a readout endpoint. Then
\[
 z_{i+1}=U_\beta z_i+aH_\beta B S_{i+1},\qquad
 Y_i=\ell z_i+\xi_i.
\]
The finite section is a stable discrete LTI system with input operator $aH_KB$.
Its transition powers have the common bound $C_*e^{-\lambda_*\Delta k}$;
\eqref{eq:finite-uniform} bounds its response sum independently of dimension.
Thus a response-level comparison cannot be blocked merely by the adjective
``infinite'' or by the boundedness of the signs.

In the cited input-Gram argument, independent isotropic sub-Gaussian inputs
supply the excitation. Its error decomposition separates response-tail,
measurement-noise and process-noise terms. Setting process noise to zero
removes those latter upper-bound terms; it does not remove input excitation.
This is a proof-level zero-process-noise specialization, not a claim that a
zero covariance satisfies a printed unit-covariance normalization. A nonzero
initial state adds a stable transient rather than an unbounded new forcing.
An applicable response estimate, plus the separately charged bias
\eqref{eq:finite-bias}, can therefore be inserted into
Theorem~\ref{thm:transfer}. An $\ell^2$ response error also bounds its maximum
norm. Physical realization by a similarity transform is unnecessary at that
interface.

What remains outside this direct regular-duration comparison is the irregular
duration mixture and history-dependent block baselines, the explicit
uniform-in-tail labeled physical inverse, and contraction of the complete
nonlinear working posterior, now with the energy budget
\eqref{eq:energy-posterior}. The auxiliary estimator above supplies a compatible
front end directly. None of these statements asserts that prior FIR methods
fail, or that a frequentist FIR bound by itself contracts an arbitrary posterior.
The cut chain is only an approximation; its removed edge is not treated as an
admissible positive infinite tail.

\section{Observable-bound rational outer certificates}
\label{sec:enclosures}
For an observable $y_q(\beta)=\sum_l\lambda_{ql}h_\beta(t_{ql})$, let the
statistical reference interval be $I_q=[\widehat y_q-R_q,\widehat y_q+R_q]$.
The rational representation must satisfy, with a separately declared
$\epsilon_q\ge0$,
\begin{equation}
 I_q\subset I_q^{\rm rat}\subset
 [\widehat y_q-R_q-\epsilon_q,\widehat y_q+R_q+\epsilon_q].
 \label{eq:representation}
\end{equation}
Rational arithmetic does not certify the coverage probability of arbitrary
supplied statistical bands. When that probability is at least $1-\alpha$,
the coverage theorem below inherits it. A nonrational band first requires a
certified rational reference enclosure, with its error explicitly charged.

Canonicalize each observable by summing rational coefficients at equal times,
dropping zero coefficients and all terms at zero time, since $h_\beta(0)=0$.
All other times must be nonnegative rational numbers. For a clock/depth
certificate, the canonical rows must be \emph{exactly} the complete set
\begin{equation}
 \bigl\{h_\beta(k\Delta+t_0)-h_\beta(k\Delta):0\le k\le N\bigr\},
 \label{eq:bound-grid}
\end{equation}
with no missing or duplicated lag. A row permutation is permitted only together
with the corresponding band permutation. Canonicalization allows algebraically
equivalent splitting of terms and arbitrary terms multiplying $h(0)$.
A count alone, or small radii alone, is not this condition.

For general rational observables, take $K\ge J+1$, mesh half-width at most
$r>0$, and Taylor order $P\ge0$. The full box variables are
$c,b_0,\ldots,b_K,a_0,\ldots,a_{K-1}$, including $a_J$. At each center,
evaluate
\[
 H_q=\sum_l\lambda_{ql}\ell\sum_{k=0}^P A_K^kB\,
                         t_{ql}^{k+1}/(k+1)!.
\]
Let $T_c$ be the largest time, $d_q=\sum_l|\lambda_{ql}|$ and
$E_0=3^{\lceil\Lambda T_c\rceil}$. Use five separate error accounts:
statistical radius, representation radius, and the following three deterministic
errors:
\begin{equation}
 \begin{split}
 E_{\rm tail}&=2T_cE_0(\Lambda T_c)^K/K!,\qquad
 E_{\rm mesh}=2rT_c^2E_0,\\
 E_{\rm num}&=T_cE_0(\Lambda T_c)^{P+1}/(P+1)!,\qquad
 E_q=d_q(E_{\rm tail}+E_{\rm mesh}+E_{\rm num}).
 \end{split}\label{eq:outer-errors}
\end{equation}
Retain a full box exactly when each $[H_q-E_q,H_q+E_q]$ intersects
$I_q^{\rm rat}$. Return both full boxes and their prefix projections.

\begin{theorem}[Bound enclosure contract]
\label{thm:outer}
The exhaustive outer union contains every originally feasible prefix.
Every admissible completion of a retained full box satisfies
\[
 |y_q(\beta)-\widehat y_q|\le r_q^{\rm out}:=R_q+\epsilon_q+2E_q.
\]
Suppose, in addition, that the model box, physical depth, target $\delta$,
clock and all recomputed certificate fields agree, that \eqref{eq:bound-grid}
is validated, and that the complete enumeration has finished. Then the returned
prefix union has diameter at most
\begin{equation}
 D_{\rm cert}=L\{2\mathcal A\max_q r_q^{\rm out}+\mathcal E\}.
 \label{eq:diameter-contract}
\end{equation}
A nonempty union with $D_{\rm cert}\le\delta$ is a certified diameter result.
An empty union is reported separately and is not a recovery assertion.
A projected prefix has an existential full-box witness; the completion claim
does not quantify over arbitrary tails attached to that projection alone.
\end{theorem}
\begin{proof}
Agreement of the full chain and cut center on boundary words through degree
$K-1$ gives the stated spatial bound. Coordinate perturbations at most $r$
change the generator by at most $4r$. Stable finite-time Duhamel and integration
give $2rT_c^2e^{\Lambda T_c}$, bounded by $E_{\rm mesh}$. Taylor remainder
gives $E_{\rm num}$, and absolute observable weights give $E_q$ uniformly over
all completions of each full box. A feasible parameter's box is retained.
At an interval intersection, one $E_q$ reaches the center and another reaches
any completion, giving $r_q^{\rm out}$. Choose full witnesses for any two
retained prefixes. The validated rows are their sampled pulse responses, so
\eqref{eq:sampled-infty} and the jet inverse give
\eqref{eq:diameter-contract}. Every step uses the bound observable identity,
not just the lengths of vectors or the magnitudes of radii.
\end{proof}

The active implementation, \texttt{tools/round53\_certificates.py}, exposes
\texttt{certify\_outer}. It recomputes its immutable model/clock certificate,
validates and orders the complete observable map with its bands, and constructs
the outer set internally. It accepts neither an external outer-set result nor
an externally supplied list of alleged outer radii. Its immutable return records
the grid, bands, model certificate, full witnesses, enumeration counts, all
five budgets and the upper bound. The arithmetic-only helper is explicitly
named \texttt{radius\_budget\_satisfied}; its return is not a grid certificate.
These statements concern outputs of the stated entry point, not authenticity
of arbitrary user-created Python objects or monkey-patched code.

Resources are checked before enumeration; exceeding them raises
\texttt{ResourceLimit}, never a successful partial certificate. Increasing
$K,P$ and then refining the mesh makes the deterministic errors small at finite
unrestricted resources. This is an exhaustive algorithm, not a claim of practical
large-depth efficiency. The statistical bands and their advertised confidence
remain a separate input premise even when a diameter certificate succeeds.

\section{All-policy depth comparison and elapsed time}
\label{sec:optimality}
Choose two admissible parameters differing only in $b_J$ by
$\delta\in(0,b_+-b_-)$, with zero initial state. Let
$q_k=(e_k,0)$, $v_k=(0,e_k)$ and $m_0=2J+1$.

\begin{lemma}[Two-way boundary propagation]
\label{lem:two-way}
For all $t\ge0$,
\[
 |g_\beta(t)-g_\eta(t)|\le
 \delta e^{\Lambda t}\Lambda^{4J+2}t^{4J+3}/(4J+3)!.
\]
The step-response bound has $t^{4J+4}/(4J+4)!$ in place of the last factor.
\end{lemma}
\begin{proof}
A generator word from $v_0$ to $q_J$ or from $v_J$ to $q_0$ needs one
velocity-to-position step and two steps per edge; its first possible degree
is $m_0$. Damping cannot shorten this. Each leg is bounded by
$e^{\Lambda t}(\Lambda t)^{m_0}/m_0!$.
The generator difference is $\pm\delta v_Jq_J^*$, so Duhamel convolves the
two legs. The beta integral
$\int_0^t(t-u)^{m_0}u^{m_0}\,du=(m_0!)^2t^{2m_0+1}/(2m_0+1)!$
gives the first inequality. Integrating in time gives the second.
\end{proof}

\begin{theorem}[All-time information bound in readout count]
\label{thm:global}
Put $\vartheta=\min\{\log2,\lambda_* /(8e\Lambda)\}$ and
$D_0=(8e\Lambda)^{-1}+C_*/\lambda_*$. For every parameter-independent
chronological policy of force at most $U$ and $n$ Gaussian boundary readouts
at selected finite times,
\[
 \KL(P_\beta^{\pi,n},P_\eta^{\pi,n})
 \le\frac{nU^2\delta^2D_0^4}{2\sigma^2}e^{-4\vartheta(2J+1)}=:K_n.
\]
Every estimator has maximal two-point probability of $d_J$ error at least
$\delta/2$ at least $\max\{0,(1-\sqrt{K_n/2})/2\}$.
\end{theorem}
\begin{proof}
Split the $L^1$ norm of each off-boundary leg at $S=m_0/(8e\Lambda)$.
The factorial short part is bounded by
$[m_0/(8e\Lambda)]4^{-m_0}\le(8e\Lambda)^{-1}2^{-m_0}$;
the stable long part by $(C_*/\lambda_*)e^{-\lambda_*S}$.
Each norm is thus at most $D_0e^{-\vartheta m_0}$.
Duhamel gives $\|g_\beta-g_\eta\|_1\le\delta D_0^2e^{-2\vartheta m_0}$.
At a common history the policy kernels coincide, and the conditional mean
difference is at most $U$ times this response norm. The chronological Gaussian
KL chain rule proves the bound. Testing the disjoint strict error balls and
using Pinsker proves the risk assertion.
\end{proof}

For fixed accuracy, $K_n\to0$ when $J_n/\log n\to\infty$, proving the
matching logarithmic depth order. With separation $4n^{-s}$ and
$J_n=\lfloor r\log n\rfloor$, the obstruction is $8\vartheta r+2s>1$;
the sufficient region is $C_\tau r+p_\tau s<1$. These are not asserted to be
the same sharp frontier. The larger lower-bound policy class and smaller
attaining committed class are compatible with this order comparison.

\begin{corollary}[Elapsed-time order with a sampling-rate restriction]
\label{cor:time}
Restrict comparison policies to interreadout gaps at least a fixed $g>0$.
At elapsed time $t$, their count is at most $1+\lfloor t/g\rfloor$.
The preceding KL bound with this count gives the logarithmic-depth obstruction
in $t$. The experiment of Theorem~\ref{thm:main}, whose stages have lengths
between $t_0$ and $t_0+T$, attains the logarithmic order under such a restriction
with $g=t_0$.
\end{corollary}
\begin{proof}
The gap bounds the number of conditional observation factors; summing the
same uniform conditional KL bound gives the claim even for a stopped record.
The deterministic two-sided stage-time bound transfers the attaining order.
\end{proof}
Without a bound on sampling rate, a fixed elapsed interval need not bound the
number of independent noisy readouts. The unrestricted policy theorem is
therefore a readout-count theorem, not an unrestricted elapsed-time theorem.

\begin{thebibliography}{99}
\bibitem{Teschl} G. Teschl, \emph{Jacobi Operators and Completely Integrable
Nonlinear Lattices}, Mathematical Surveys and Monographs 72, AMS, 2000.
\bibitem{MM2019} A. S. Mikhaylov and V. S. Mikhaylov,
Dynamic inverse problem for Jacobi matrices,
\emph{Inverse Problems and Imaging} 13 (2019), 431--447.
\href{https://doi.org/10.3934/ipi.2019021}{doi:10.3934/ipi.2019021}.
\bibitem{SRD2020} T. Sarkar, A. Rakhlin and M. A. Dahleh,
Nonparametric Finite Time LTI System Identification,
\href{https://arxiv.org/abs/1902.01848v6}{arXiv:1902.01848v6}, 8 April 2020.
Section numbering in the comparison refers to this preprint.
Journal version: \emph{Finite Time LTI System Identification},
JMLR 22(26) (2021), 1--61.
\bibitem{SS2026} F. Seizilles and M. Siebel,
Posterior contraction under misspecification and heteroscedasticity in
non-linear inverse problems,
\href{https://arxiv.org/abs/2603.28177v1}{arXiv:2603.28177v1}, 30 March 2026.
\end{thebibliography}
\end{document}
